\paragraph{Step 22 finite-carrier F51 statement.}
Let \(H\) be the complete finite mode carrier over
\((d0,d1,d2,d3)\).  Let
\[
q_U=\pi_L:H\to O_L,\qquad
q_A=q_{\mathrm{QM}}:H\to O_{\mathrm{QM}},\qquad
q_B=q_{\mathrm{GR}}:H\to O_{\mathrm{GR}}.
\]
The parent projections are
\[
\pi_A:O_L\to O_{\mathrm{QM}},\qquad
\pi_B:O_L\to O_{\mathrm{GR}}.
\]
The F51 object equations are checked by
\[
\pi_A\circ q_U=q_A,\qquad \pi_B\circ q_U=q_B.
\]
On the Step 22 carrier both children have nonempty factorization
defect against \(q_U\), so \(q_U\) is a strict common refinement of
both child quotients.

\paragraph{Status compatibility.}
The parent audit \(A_U=(0.70,0.20,-0.40,0.35)\) restricts to
\[
A_{\mathrm{QM}}=(0.70,0.20,-0.40),\qquad
A_{\mathrm{GR}}=(0.70,-0.40,0.35).
\]
On the shared overlap \((d0,d2)\), the two child audits agree:
\[
(0.70,-0.40)=(0.70,-0.40).
\]
Thus the lifted claim/status records are compatible on the finite
carrier.  The status-conflict control changes the GR overlap prices,
making this equality fail, and therefore fails F51 status
compatibility.

\paragraph{Verdict.}
On the Step 22 carrier, \(L\) satisfies the F51 common-refinement
unification criterion for \(q_{\mathrm{QM}}\) and \(q_{\mathrm{GR}}\).
